El auge de las plataformas de \emph{streaming} ha incrementado la saturación en la industria de la música, lo que impulsa a artistas a buscar nuevas formas de diferenciación y conexión con sus audiencias. En este contexto, la realidad virtual (VR por sus siglas en inglés) surge como una alternativa para crear experiencias musicales inmersivas, interactivas y personalizadas. Sin embargo, gran parte de los estudios sobre experiencias musicales en VR se orientan a conciertos en vivo, navegación activa, múltiples perspectivas o interacción social entre audiencia e intérprete. Esto deja una oportunidad para explorar experiencias tipo videoclip musical, donde los elementos visuales, auditivos e interactivos sean organizados mediante una secuenciación temporal diseñada para acompañar la narrativa e identidad de una canción específica.

Este proyecto propone el desarrollo de un sistema de secuenciación inmersiva basado en Arquitecturas Dirigidas por Eventos (EDA por sus siglas en inglés) que permite organizar y coordinar elementos visuales, interactivos y auditivos dentro de un videoclip musical en VR. A través de eventos temporizados, se sincronizan cambios de ambiente, estímulos visuales, interacciones y elementos sonoros con la estructura de la canción, de manera que la experiencia mantenga coherencia con su narrativa musical.

La experiencia es evaluada mediante el \emph{I-Group Presence Questionnaire} (IPQ), instrumento que permite medir la percepción de presencia mediante dimensiones como la presencia espacial, el involucramiento y el realismo. Con ellos, se busca analizar si una secuenciación basada en eventos contribuye a fortalecer la percepción de presencia del usuario y ofrecer una experiencia VR personalizada, inmersiva y alineada con la intención artística del proyecto.
